%% ma: L12-B16-0010
%% lop: 12
%% chuong: 5
%% bai: 16
%% dang: 12.16.8
%% loai: TN4
%% muc: NB
%% dap_an: "A"
%% nguon: "Ảnh giáo viên gửi 04/11/2026 (Huỳnh Văn Ánh, Chuyên đề V - Phương pháp tọa độ trong không gian), Đề test theo bài, Phần I Câu 1"
%% kien_thuc: "Bài 16 (góc giữa hai đường thẳng), Bài 8 (tích vô hướng, độ dài vectơ)"
%% trang_thai: cho-duyet
%% ly_do: "Dạng 12.16.8 mới đề xuất, chờ giáo viên duyệt."
%% ghi_chu: "Đáp án gốc A khớp (sympy)."
\begin{ex}
Trong không gian với hệ tọa độ $Oxyz$, cho hai đường thẳng $d_1$ và $d_2$ lần lượt có vectơ chỉ phương là $\vec u_1=(1;-2;-3)$ và $\vec u_2=(-4;1;5)$. Góc giữa hai đường thẳng $d_1,d_2$ là
\choice{\True $30^\circ$}{$45^\circ$}{$90^\circ$}{$60^\circ$}
\loigiai{$\cos(d_1,d_2)=\dfrac{|\vec u_1\cdot\vec u_2|}{|\vec u_1|\cdot|\vec u_2|}=\dfrac{|-4-2-15|}{\sqrt{14}\cdot\sqrt{42}}=\dfrac{21}{14\sqrt3}=\dfrac{\sqrt3}{2}$. Suy ra $(d_1,d_2)=30^\circ$. Chọn A.}
\end{ex}
